% Library unit: the rule of the square. Printed once, before the first position that uses it; see lib/README.md.

\paragraph{The rule of the square}\label{rule-of-the-square}\index{King and pawn endings!Rule of the square}
When the attacking king is far away, a race between the pawn and the defending king can be
decided without counting moves. Draw a square from the pawn to its queening square, as in the
diagrams below. This is the \termas{rule-of-the-square}{rule of the
square}: if the defending king can step into the
square, it catches the pawn; if it cannot, the pawn queens.

\keysquares{8/8/8/7k/2P5/8/8/7K b - - 0 1}{c4,c5,c6,c7,c8,d4,d5,d6,d7,d8,e4,e5,e6,e7,e8,f4,f5,f6,f7,f8,g4,g5,g6,g7,g8}{Black to move: 1...\K g5 steps
into the square and catches the pawn. A draw.}

\keysquares{8/8/8/2P4k/8/8/8/7K b - - 0 1}{c5,c6,c7,c8,d5,d6,d7,d8,e5,e6,e7,e8,f5,f6,f7,f8}{White to move instead plays
1.c5. The square shrinks, the black king cannot get in, and the pawn queens.}

Two things to watch. A pawn on its starting square can move two squares at once, so draw its
square as if it stood one square further. And the rule only counts moves: it does not see a king
or a pawn standing in the way of the defending king.
